% block0_pde_foundations.tex
%
% Phase 3, Block 0: Foundations of the finite-difference solution of the
% Black-Scholes PDE.
%
% Self-contained writeup. Compiles standalone with pdflatex.

\documentclass[11pt,a4paper]{article}

\usepackage[T1]{fontenc}
\usepackage[utf8]{inputenc}
\usepackage{amsmath, amssymb, amsthm}
\usepackage{geometry}
\usepackage{hyperref}
\usepackage{enumitem}

\geometry{margin=1in}

\newtheorem{proposition}{Proposition}[section]
\newtheorem{theorem}[proposition]{Theorem}
\newtheorem{definition}[proposition]{Definition}
\newtheorem{remark}[proposition]{Remark}

\title{Phase 3, Block 0:\\Foundations of the finite-difference solution\\of the Black-Scholes PDE}
\author{Quant Finance Portfolio}
\date{}

\begin{document}
\maketitle

\section{Scope}

This block sets up the mathematical and numerical scaffolding on which
all of Phase~3 rests. It contains no pricer. Its deliverables are:
\begin{enumerate}[leftmargin=2em]
\item the logarithmic transformation that turns the Black--Scholes PDE
      into a constant-coefficient parabolic PDE;
\item the truncation of the spatial domain and the choice of artificial
      Dirichlet boundary conditions;
\item the discretisation grid and the standard centred finite-difference
      stencils;
\item the von~Neumann stability analysis applied to the explicit (FTCS)
      and implicit (BTCS) schemes for the pure-diffusive part of the PDE.
\end{enumerate}
The pricers themselves (FTCS, BTCS, Crank--Nicolson, PSOR, trinomial)
appear in Blocks~1--5, all built on this foundation.

% =====================================================================
\section{The Black--Scholes PDE}
% =====================================================================

Under the risk-neutral measure $\mathbb{Q}$, the price $V(S,t)$ of a
European option on a non-dividend-paying asset $S_t$ following geometric
Brownian motion
\[
dS_t = r S_t\, dt + \sigma S_t\, dW_t^{\mathbb{Q}}
\]
satisfies the Black--Scholes partial differential equation
\begin{equation}
\frac{\partial V}{\partial t}
+ \frac{1}{2}\sigma^2 S^2 \frac{\partial^2 V}{\partial S^2}
+ r S \frac{\partial V}{\partial S}
- r V \;=\; 0,
\qquad S > 0,\; t \in [0,T),
\label{eq:BS}
\end{equation}
with terminal condition $V(S,T) = \phi(S)$, where $\phi$ is the option
payoff. This is a backward parabolic equation in calendar time.

Direct discretisation of \eqref{eq:BS} on a uniform grid in $S$ is
unsatisfactory for two independent reasons. First, the coefficients
$\frac{1}{2}\sigma^2 S^2$ and $rS$ depend on the grid position, which
complicates the von~Neumann stability analysis: the standard analysis
assumes constant coefficients. Second, a uniform grid in $S$ resolves
the lognormal-tailed distribution of $S_T$ poorly, with too few nodes
near $S = 0$ relative to the scale on which $V(S,t)$ varies.

Both problems are removed by a single change of variables.

% =====================================================================
\section{Logarithmic transformation}\label{sec:logtransform}
% =====================================================================

Introduce
\begin{equation}
x \;=\; \ln\!\left(\frac{S}{K}\right), \qquad
\tau \;=\; T - t.
\label{eq:transform}
\end{equation}
The variable $x$ is the log-moneyness; $\tau \in [0, T]$ is the time to
maturity. Anchoring at the strike $K$ (rather than using $\ln S$
directly) places $x = 0$ at the at-the-money point, which gives a
canonical centring for the truncated domain.

Computing the derivatives by the chain rule, with $S = K e^x$,
\[
\frac{\partial V}{\partial t} \;=\; -\frac{\partial V}{\partial \tau},
\qquad
\frac{\partial V}{\partial S} \;=\; \frac{1}{S}\frac{\partial V}{\partial x},
\qquad
\frac{\partial^2 V}{\partial S^2}
\;=\; \frac{1}{S^2}\!\left(
  \frac{\partial^2 V}{\partial x^2} - \frac{\partial V}{\partial x}
\right).
\]
Substituting into \eqref{eq:BS} and rearranging yields
\begin{equation}
\boxed{\;
\frac{\partial V}{\partial \tau}
\;=\; \frac{\sigma^2}{2}\frac{\partial^2 V}{\partial x^2}
   + \mu\,\frac{\partial V}{\partial x}
   - r\, V,
\qquad
\mu \;:=\; r - \frac{\sigma^2}{2}.
\;}
\label{eq:transformed}
\end{equation}
This is a forward parabolic equation in $\tau$ on the unbounded domain
$x \in \mathbb{R}$, with constant coefficients. The drift $\mu$ in
\eqref{eq:transformed} is precisely the drift of $\ln S_t$ under
$\mathbb{Q}$, by It\^o.

The initial condition becomes
\[
V(x, 0) \;=\; \phi(K e^x), \qquad x \in \mathbb{R}.
\]

\begin{remark}
The transformation \eqref{eq:transform} preserves the regularity of
$V$. The non-smoothness of $\phi$ at the strike (the kink of
$(\cdot)^+$) becomes a kink of the transformed initial condition at
$x = 0$. This irregularity will be relevant in Block~3 when we discuss
spurious oscillations in Crank--Nicolson.
\end{remark}

% =====================================================================
\section{Structure of the transformed PDE}
% =====================================================================

Equation \eqref{eq:transformed} is a constant-coefficient
\emph{convection--diffusion--reaction} (CDR) equation:
\[
V_\tau \;=\; D\, V_{xx} \,+\, a\, V_x \,+\, c\, V,
\qquad
D = \frac{\sigma^2}{2} > 0,\quad
a = \mu \in \mathbb{R},\quad
c = -r \le 0.
\]
Each of the three terms has a distinct numerical character.
\begin{itemize}[leftmargin=2em]
\item The \emph{diffusive} term $D\, V_{xx}$ is dominant for typical
      financial parameters ($\sigma \in [0.1, 0.5]$, $\Delta x$ small).
      It governs the most restrictive stability constraint on the
      explicit scheme of Block~1.
\item The \emph{convective} term $a\, V_x$ encodes directional drift.
      For $\mu > 0$ information flows toward $+x$; for $\mu < 0$
      toward $-x$. With centred space differences (used throughout
      Phase~3) the contribution to the von~Neumann factor is purely
      imaginary.
\item The \emph{reactive} term $c V = -rV$ is local in space and acts
      as a uniform exponential decay with rate $r$. It contributes
      only a multiplicative factor $1 - r\Delta\tau$ (or its implicit
      analogue) to the amplification factor and never destabilises
      the scheme.
\end{itemize}

\begin{remark}[Why we keep all three terms together]
A common simplification in textbooks is to apply a further change of
variables $V = e^{-r\tau} U$ (kills the reactive term) and a Galilean
shift in $x$ (kills the convective term), reducing
\eqref{eq:transformed} to the pure heat equation $U_\tau = D\,U_{xx}$.
This is mathematically clean but obscures the discrete structure: the
boundary conditions and the payoff become time-dependent in non-trivial
ways, and the relationship to the original $V$ is harder to read off
the code. Following the standard practice of Tavella \&
Randall~\cite{TavellaRandall2000}, we discretise
\eqref{eq:transformed} as it stands.
\end{remark}

% =====================================================================
\section{Truncation of the spatial domain}
% =====================================================================

The unbounded domain $x \in \mathbb{R}$ has to be truncated to a finite
interval $[x_{\min}, x_{\max}]$. The truncation introduces an artificial
boundary; on it we will impose Dirichlet conditions whose values are
read off the analytic asymptotics of $V$.

A standard choice is symmetric truncation centred at $x = 0$:
\begin{equation}
x_{\min} \;=\; -N_\sigma\,\sigma\sqrt{T},
\qquad
x_{\max} \;=\; +N_\sigma\,\sigma\sqrt{T},
\label{eq:domain}
\end{equation}
with $N_\sigma$ a small constant (typically $4$ or $5$). The
justification is that the lognormal density of $S_T = K e^{x_T}$ has
$x_T \sim \mathcal{N}(\mu T,\, \sigma^2 T)$ under
$\mathbb{Q}$, so the probability mass outside
$[-N_\sigma \sigma\sqrt{T},\, +N_\sigma \sigma\sqrt{T}]$ is
$O\!\bigl(e^{-N_\sigma^2/2}\bigr)$, which for $N_\sigma = 4$ is
$\approx 6 \times 10^{-5}$ and decays super-polynomially in $N_\sigma$.

\begin{remark}
The total error of any FD pricer is the sum of three contributions:
the discretisation error (powers of $\Delta x$ and $\Delta\tau$), the
truncation error from the artificial boundary, and floating-point
round-off. For $N_\sigma \ge 4$ the truncation contribution is
typically below the discretisation error already at moderate
$N \sim 100$ and never becomes the dominant error, so we treat it as
negligible throughout. This separation of error sources will reappear
in Block~3 when verifying the empirical convergence rate of CN.
\end{remark}

% =====================================================================
\section{Boundary and initial conditions}\label{sec:bcs}
% =====================================================================

\subsection*{European call}

The terminal payoff $\phi(S) = (S-K)^+$ becomes, in $x$,
\begin{equation}
V(x, 0) \;=\; K\,(e^x - 1)^+,
\qquad x \in [x_{\min}, x_{\max}].
\label{eq:ic-call}
\end{equation}
For the boundary conditions we use the well-known asymptotic limits.
As $S \to 0$ the call goes to zero; as $S \to \infty$ it tends to the
forward $S - K e^{-r(T-t)}$. Translating to $(x, \tau)$,
\begin{align}
V(x_{\min}, \tau) &\;=\; 0, \label{eq:bc-call-lo}\\
V(x_{\max}, \tau) &\;=\; K\!\left(e^{x_{\max}} - e^{-r\tau}\right).
\label{eq:bc-call-hi}
\end{align}
Both equalities hold in the limit $x_{\min}\to -\infty$,
$x_{\max} \to +\infty$. On a truncated domain they are imposed as
exact Dirichlet conditions; the small additional error this introduces
is the truncation error discussed above.

\subsection*{European put}

By analogous reasoning, the put payoff $\phi(S) = (K-S)^+$ becomes
\begin{equation}
V(x, 0) \;=\; K\,(1 - e^x)^+,
\label{eq:ic-put}
\end{equation}
and the asymptotic boundary values are
\begin{align}
V(x_{\min}, \tau) &\;=\; K\!\left(e^{-r\tau} - e^{x_{\min}}\right),
   \label{eq:bc-put-lo}\\
V(x_{\max}, \tau) &\;=\; 0.
   \label{eq:bc-put-hi}
\end{align}
Equation \eqref{eq:bc-put-lo} is a simple consequence of the
martingale property: as $S \to 0$, $\mathbb{Q}\{S_T < K\} \to 1$ so
\(
V(S,t) \to e^{-r(T-t)}\bigl(K - \mathbb{E}^{\mathbb{Q}}[S_T \mid S_t = S]\bigr)
= K e^{-r(T-t)} - S,
\)
which is precisely $K(e^{-r\tau} - e^x)$.

\begin{remark}
For Block~4 (American put) the boundary conditions
\eqref{eq:bc-put-lo}--\eqref{eq:bc-put-hi} have to be replaced by the
\emph{intrinsic-value} conditions $V(x_{\min},\tau) = K - K e^{x_{\min}}$
(the option is exercised immediately for very low $S$) and the same
$V(x_{\max},\tau) = 0$. We will revisit this in Block~4.
\end{remark}

% =====================================================================
\section{The discretisation grid}
% =====================================================================

We discretise on a uniform tensor-product grid in $(x, \tau)$. Let
\begin{equation}
N \in \mathbb{N},\quad M \in \mathbb{N},\qquad
\Delta x \;=\; \frac{x_{\max} - x_{\min}}{N},\qquad
\Delta \tau \;=\; \frac{T}{M}.
\label{eq:grid-spacing}
\end{equation}
The spatial nodes are
\(
x_j = x_{\min} + j\,\Delta x,\; j = 0, 1, \dots, N,
\)
giving $N+1$ nodes; the time levels are
\(
\tau_n = n\,\Delta \tau,\; n = 0, 1, \dots, M,
\)
giving $M+1$ levels.

The discrete unknown is $V_j^n \approx V(x_j, \tau_n)$. The
\emph{interior nodes} are $j = 1, \dots, N-1$ (a total of $N-1$
unknowns at each time level); the boundary nodes $j = 0$ and $j = N$
are fixed by the Dirichlet conditions of Section~\ref{sec:bcs}. At
$\tau_0 = 0$ all nodes are fixed by the initial condition; the answer
of interest is $V(0, T)$ if $S_0 = K$, or more generally the
interpolant of $V_j^M$ at $x = \ln(S_0/K)$.

The uniformity of the grid in $x$ (not in $S$) is the consequence of
the logarithmic transformation: a uniform grid in $S$ would
over-resolve the region near $S = 0$ and under-resolve large $S$.
The transformation makes the natural sampling uniform in
log-moneyness, which is also the natural scale for $V$ to vary on.

% =====================================================================
\section{Finite-difference stencils}
% =====================================================================

Let $f$ be a smooth function. Taylor expansion gives
\[
f(x \pm h) \;=\; f(x) \pm h f'(x) + \frac{h^2}{2} f''(x)
   \pm \frac{h^3}{6} f'''(x) + O(h^4).
\]
Combining the $+h$ and $-h$ expansions yields the centred stencils
\begin{align}
f'(x)  &\;\approx\; \frac{f(x+h) - f(x-h)}{2h} + O(h^2),
   \label{eq:fd-d1}\\
f''(x) &\;\approx\; \frac{f(x+h) - 2 f(x) + f(x-h)}{h^2} + O(h^2).
   \label{eq:fd-d2}
\end{align}
For the time derivative we use forward or backward Euler:
\begin{align}
\partial_\tau V|_j^n &\;\approx\;
   \frac{V_j^{n+1} - V_j^n}{\Delta \tau} + O(\Delta \tau),
   \qquad &\text{(forward, used in FTCS)}\label{eq:fd-fwd}\\
\partial_\tau V|_j^n &\;\approx\;
   \frac{V_j^n - V_j^{n-1}}{\Delta \tau} + O(\Delta \tau).
   \qquad &\text{(backward, used in BTCS)}\label{eq:fd-bwd}
\end{align}
Centred-in-time differences (used in Crank--Nicolson, Block~3) achieve
$O(\Delta\tau^2)$ at the cost of being implicit at both adjacent
levels.

The label \emph{FTCS} stands for \emph{Forward-Time, Centred-Space}
(combining \eqref{eq:fd-fwd} with \eqref{eq:fd-d1}--\eqref{eq:fd-d2}
evaluated at level $n$); \emph{BTCS} for \emph{Backward-Time,
Centred-Space} (combining \eqref{eq:fd-bwd} with the same spatial
stencils, but evaluated at level $n+1$ from the perspective of the
unknown).

% =====================================================================
\section{Von Neumann stability analysis}\label{sec:vN}
% =====================================================================

Stability of a finite-difference scheme is the property that
round-off errors do not grow unboundedly as the time-stepping
proceeds. Together with consistency (the truncation error tends to
zero as $\Delta x, \Delta\tau \to 0$), it implies convergence:
\begin{theorem}[Lax equivalence,
{\cite[Theorem 1.5.1]{Strikwerda2004}}]
For a well-posed linear initial-value problem and a consistent
finite-difference approximation, stability is necessary and sufficient
for convergence.
\end{theorem}
The criterion of stability we will use is the von~Neumann criterion,
which exploits the fact that on a uniform grid with constant
coefficients the discrete Fourier modes $e^{ij\theta}$, $\theta \in
[-\pi, \pi]$, diagonalise the spatial operators.

\subsection{Setup}

Write the round-off error at the level $n$ as a single Fourier mode,
\begin{equation}
\varepsilon_j^n \;=\; g(\theta)^n\, e^{ij\theta},
\label{eq:vn-ansatz}
\end{equation}
with $g(\theta) \in \mathbb{C}$ the (mode-dependent) \emph{amplification
factor}. Linearity of the scheme allows independent treatment of each
$\theta$. The scheme is \emph{von Neumann stable} iff
\begin{equation}
|g(\theta)| \;\le\; 1
\qquad
\text{for all }\theta \in [-\pi, \pi].
\label{eq:vn-criterion}
\end{equation}

\subsection{FTCS for the pure-diffusive equation}

Consider first the model problem $V_\tau = D V_{xx}$, ignoring the
convective and reactive terms; we will reinstate them below. The FTCS
discretisation is
\[
\frac{V_j^{n+1} - V_j^n}{\Delta \tau}
\;=\; D\,\frac{V_{j+1}^n - 2 V_j^n + V_{j-1}^n}{\Delta x^2}.
\]
Substituting the ansatz \eqref{eq:vn-ansatz}, dividing by
$g^n e^{ij\theta}$, and using $e^{i\theta} + e^{-i\theta} = 2\cos\theta$:
\[
g - 1 \;=\; \frac{D\,\Delta\tau}{\Delta x^2}\,
   \bigl(2\cos\theta - 2\bigr)
\;=\; -4\,\alpha\,\sin^2\!\tfrac{\theta}{2},
\]
where the dimensionless \emph{Fourier number} is
\begin{equation}
\alpha \;:=\; D\,\frac{\Delta\tau}{\Delta x^2}
       \;=\; \frac{\sigma^2}{2}\cdot\frac{\Delta\tau}{\Delta x^2}.
\label{eq:fourier-number}
\end{equation}
Therefore
\[
g(\theta) \;=\; 1 - 4\alpha\,\sin^2\!\tfrac{\theta}{2}\;\in\;\mathbb{R}.
\]
The criterion \eqref{eq:vn-criterion} becomes $-1 \le g \le 1$. The
upper bound is automatic since $\alpha,\sin^2 \ge 0$. The lower bound
is tight at $\theta = \pi$ and gives
\begin{equation}
\boxed{\;\alpha \le \tfrac{1}{2}\;}\quad \iff \quad
\frac{\sigma^2 \Delta\tau}{\Delta x^2} \le 1.
\label{eq:cfl}
\end{equation}
\eqref{eq:cfl} is the \emph{CFL condition} for the heat equation, in
honour of Courant, Friedrichs and Lewy. It says that an explicit
diffusive scheme cannot propagate the heat by more than half a grid
spacing per time step. The constant $1/2$ is sharp.

\begin{remark}
The CFL constraint is severe in practice. For $\sigma = 0.2$,
$\Delta x = 0.05$, the largest stable $\Delta\tau$ is
$\alpha\Delta x^2 / D \approx \frac{0.5 \cdot 0.0025}{0.02} = 0.0625$,
so a year requires at least $T/\Delta\tau = 16$ time steps. Halving
$\Delta x$ requires \emph{quadrupling} the number of time steps. This
$O(\Delta x^{-2})$ scaling of the time-step count is the principal
limitation of FTCS and the motivation for the implicit schemes in
Block~2 and beyond.
\end{remark}

\subsection{BTCS for the pure-diffusive equation}

The implicit scheme evaluates the spatial stencil at level $n+1$:
\[
\frac{V_j^{n+1} - V_j^n}{\Delta \tau}
\;=\; D\,\frac{V_{j+1}^{n+1} - 2 V_j^{n+1} + V_{j-1}^{n+1}}{\Delta x^2}.
\]
Substituting \eqref{eq:vn-ansatz} gives
\[
g - 1 \;=\; -4\alpha\,g\,\sin^2\!\tfrac{\theta}{2}
\;\Longrightarrow\;
g(\theta) \;=\; \frac{1}{1 + 4\alpha\sin^2(\theta/2)}.
\]
Since $\alpha > 0$, the denominator is at least $1$ for every
$\theta$, so $g \in (0, 1]$. The criterion is satisfied for any
$\alpha > 0$:
\begin{proposition}[Unconditional stability of BTCS]
\label{prop:btcs}
The BTCS scheme for $V_\tau = D V_{xx}$ is von Neumann stable for any
$\Delta\tau > 0,\; \Delta x > 0$.
\end{proposition}

This is the central trade-off of Phase~3. FTCS is cheap per step
(matrix-free) but has the CFL constraint; BTCS allows arbitrary
$\Delta\tau$ but requires solving a linear system at each step.
Block~2 will pay the cost of that linear solve via the Thomas
algorithm.

\subsection{Convective term: an aside}

Reinstating the convective term $a V_x$ with centred-space
differences in FTCS, the amplification factor becomes complex,
\[
g(\theta) \;=\; 1 - 4\alpha\sin^2\!\tfrac{\theta}{2}
   + i\,\nu\,\sin\theta,
\qquad
\nu \;:=\; a\,\frac{\Delta\tau}{2\,\Delta x},
\]
where $\nu$ is the \emph{Courant number} (advective). Imposing
$|g|^2 \le 1$ yields a mixed condition coupling $\alpha$ and $\nu$;
algebraically, after some manipulation (see~\cite[\S3.5]{Strikwerda2004}),
\[
4\alpha^2 \;\ge\; \nu^2
\quad\text{and}\quad
\alpha \le \tfrac{1}{2}.
\]
For Black--Scholes parameters the relevant figure of merit is the
\emph{grid Péclet number} $\mathrm{Pe}_h = |a|\Delta x / (2D) =
|\mu|\Delta x / \sigma^2$. For $\sigma = 0.2$, $|\mu| \le 0.1$,
$\Delta x \le 0.1$, $\mathrm{Pe}_h \le 0.25$, so the convective
contribution is sub-dominant and the diffusive bound \eqref{eq:cfl}
governs in practice. The reactive term $-rV$, being local in space,
contributes only a multiplicative factor $1 - r\Delta\tau$ to $g$ and
does not affect stability for $\Delta\tau \le 1/r$, which is always
satisfied.

\subsection{Summary: two control numbers}

The discrete dynamics of the entire phase is controlled by two
dimensionless numbers,
\begin{equation}
\alpha \;=\; \frac{\sigma^2}{2}\cdot\frac{\Delta\tau}{\Delta x^2},
\qquad
\nu \;=\; \frac{|\mu|}{2}\cdot\frac{\Delta\tau}{\Delta x}.
\label{eq:control-numbers}
\end{equation}
FTCS (Block~1) is stable iff $\alpha \le 1/2$ and the convective
contribution does not violate the mixed bound; in the parameter range
of this project this reduces effectively to $\alpha \le 1/2$.
BTCS (Block~2) and Crank--Nicolson (Block~3) are unconditionally stable
in $\alpha$. As we will see in Block~5, the Kamrad--Ritchken trinomial
tree maps to FTCS under a specific reparametrisation; its
non-negative-probability condition $\lambda^2 \ge 1/\alpha$ is the
discrete-tree avatar of the CFL constraint \eqref{eq:cfl}.

% =====================================================================
\section{Implementation outline}
% =====================================================================

The Phase 3 Block 0 implementation provides four ingredients:
\begin{enumerate}[leftmargin=2em]
\item \texttt{Grid} (Python dataclass / C++ struct): an immutable
      container holding $N$, $M$, $T$, $\sigma$, $r$, $K$,
      $x_{\min}$, $x_{\max}$, $\Delta x$, $\Delta \tau$, the arrays
      $\{x_j\}$ and $\{\tau_n\}$.
\item \texttt{build\_grid}: a constructor that picks
      $x_{\min}, x_{\max}$ from
      $\pm N_\sigma \sigma \sqrt{T}$ (default $N_\sigma = 4$) and
      assembles the grid.
\item Initial-condition functions
      \texttt{call\_initial\_condition}, \texttt{put\_initial\_condition}
      returning the payoff vector at $\tau = 0$.
\item Boundary functions \texttt{call\_boundary\_lower},
      \texttt{call\_boundary\_upper}, \texttt{put\_boundary\_lower},
      \texttt{put\_boundary\_upper} returning Dirichlet values for
      every $\tau_n$.
\item Stability-number utilities \texttt{fourier\_number} and
      \texttt{courant\_number} for diagnostic use by the explicit
      pricer in Block~1.
\end{enumerate}
This module imports nothing from outside \texttt{numpy} and the
standard library. It contains no pricer.

% =====================================================================
\section*{References}\label{sec:refs}
% =====================================================================

\begin{thebibliography}{9}

\bibitem{TavellaRandall2000}
D. Tavella and C. Randall.
\emph{Pricing Financial Instruments: The Finite Difference Method.}
Wiley, 2000. Chapters 2--4.\\
{\small The reference for finite differences in finance: log
transformation, domain truncation, boundary conditions, and stability
analysis specific to Black--Scholes.}

\bibitem{Seydel2017}
R. Seydel.
\emph{Tools for Computational Finance.}
6th ed., Springer, 2017. Chapter~4.\\
{\small Compact, implementation-oriented treatment.}

\bibitem{MortonMayers2005}
K.~W. Morton and D.~F. Mayers.
\emph{Numerical Solution of Partial Differential Equations.}
2nd ed., Cambridge University Press, 2005. Chapters~2--3.\\
{\small Rigorous mathematical treatment of consistency, stability and
the Lax equivalence theorem.}

\bibitem{Strikwerda2004}
J.~C. Strikwerda.
\emph{Finite Difference Schemes and Partial Differential Equations.}
2nd ed., SIAM, 2004. Chapters~1--3.\\
{\small Authoritative reference for von Neumann analysis applied to
linear constant-coefficient schemes.}

\bibitem{Wilmott1998}
P. Wilmott.
\emph{Derivatives: The Theory and Practice of Financial Engineering.}
Wiley, 1998. Chapters~46--49.\\
{\small Practitioner-oriented exposition with explicit Black--Scholes
PDE discretisation.}

\end{thebibliography}

\end{document}
